% 第 16 章　为什么世界总爱振动
\chapter{为什么世界总爱振动}

\step{反常识开场}

先看三件你多半见过、却很少追问的小事。

第一件在秋千上。一个人站在秋千旁边，不用大力气，只在秋千每次荡到最高点的瞬间轻轻推一下——力气小到几乎像碰一碰——几分钟之后，秋千竟然荡得比谁都高。如果他用大得多的力气乱推一气，有时推在秋千回来的时候，秋千反而很快慢下来。

第二件在厨房。把一根不锈钢饭勺挂在手指上，让它像钟摆一样左右晃，再用手机计时。你会隐约发现一个奇怪的事实：不管它晃得宽一点还是窄一点，来回一次花的时间几乎一样。晃宽了它似乎“急”一些，晃窄了它似乎“懒”一些，可周期就是不配合你的直觉。

第三件在桥上。1940 年，美国塔科马海峡大桥在一场并不算罕见的大风中像麻花一样扭动，最后折断坠入海峡，全程被拍成了电影胶片。一座钢结构大桥，怎么会像琴弦一样“振”起来？

这三件事背后是同一条线索：这个世界上的东西，一旦偏离了它“想待”的位置，往往不是安静地回去，而是要来来回回折腾好一阵。吉他弦被拨一下会响，钟摆被碰一下会摆，汽车碾过减速带会颠两三下才稳，杯子里的水面被磕一下会晃，甚至连晶格里的原子、麦克风里的电路、核磁共振仪里的原子核，也都在以各自的节奏振动。为什么世界这么爱振动？这一章就来回答这个问题，顺便解释：为什么轻轻一推，也能把秋千推上天。

值得先强调一句：这些现象你全都见过，本章不会向你展示任何你接触不到的新事实。我们要做的，是把这些散落的事实串进同一个框架，让你下次看见晾衣绳在风里抖、听见冰箱压缩机嗡嗡响、感觉地铁启动时吊环来回晃的时候，能脱口而出：哦，又是一个振子。

\step{先猜一猜}

在继续读之前，请把你的直觉押在下面几个问题上，不许反悔。

\begin{enumerate}[leftmargin=2em]
  \item 秋千上换一个体重加倍的人，荡一个来回的时间会怎么变：明显变长、明显变短，还是几乎不变？
  \item 想让秋千越荡越高，关键在于推力大，还是推得“是时候”？
  \item 一根弹簧下面挂一个重物，拉下来放手后上下振动。如果把这套装置搬到月球上去（重力只有地球上的六分之一），它振动的快慢会变吗？
  \item 汽车碾过一个坑之后为什么总要颠几下才平稳，而不是只沉一下就停住？如果去掉减震器，会颠得更多还是更少？
\end{enumerate}

把答案写下来。这一章结束时，你再回头批改自己的卷子。多数人至少会答错一道，而答错的那一道，往往正是本章的核心。先透露一点：这四道题的答案都不依赖任何新定律，只靠牛顿定律加上一个本章要发明的新视角就能推出来。物理学的进步常常不是发现新现象，而是学会用新的方式整理旧现象——振动就是这样一种整理方式。

\step{动手实验}

\begin{experiment}[桌边的钢尺和弹簧上的钥匙串]
实验一：把一把钢尺（或塑料尺）压在桌边，伸出一截在桌面外。用一只手死死按住桌面上的部分，另一只手把伸出的部分向下扳一点再松手。尺子会嗡嗡地振动并发出声音。逐次改变伸出桌外的长度，你会发现：伸出越短，振动越快，音调越高；伸出越长，振动越慢，音调越低。一个不用任何仪器就能得到的结论是：振动有它自己的节奏，而且这个节奏由“系统本身”决定——尺子、伸出长度，而不是由你扳得多狠决定。你扳得再狠，它也只是振得更“用力”，节奏基本不变。

实验二：找一根橡皮筋或细弹簧，下面挂一串钥匙。提着手让钥匙串静止，然后轻轻向下一拉再松手。钥匙串会上下振动很久。数一数十秒内振动了多少次。接着把钥匙串减半（只留一半钥匙），再数一次。多数人的直觉是“轻了振得慢”或者“轻了振得快”，请用你自己的数据回答。

实验三：让钥匙串静止下垂，你的手做极小幅度的上下抖动。先胡乱地快抖，钥匙串没什么反应；再把抖动的节奏放慢，慢慢逼近你刚才数出的那个“固有节奏”。当节奏对上的一刹那，钥匙串会突然像疯了一样大幅度上下跳动，而你的手动得比刚才还小。请记住这个感觉：对上节奏，比使大力气重要得多。如果你愿意再进一步，试着在对上节奏之后故意把抖动放慢一点点，你会发现钥匙串的大幅跳动维持不住，慢慢萎下去——供能的“存折”只在特定节奏上开门。
\end{experiment}

这三个实验都不需要任何公式，却已经包含了本章的全部关键词：平衡位置、固有节奏、振幅、对上节奏。下面我们给它们正式命名。

如果做实验时有条件，建议再用手机慢动作录像回放钥匙串的运动，特别注意两件事：一是钥匙串经过最低点附近时看起来“一晃而过”，在两端的最高点和最低点却“磨蹭”很久——说明速度在中间最大、在两端为零；二是不管振幅大小，数出来的节奏几乎不变。这两个观察不用任何理论就能得到，却正是后面公式的全部经验基础。一个好理论不该发明事实，只该解释你亲手看到的事实。

\step{旧模型遇到麻烦}

我们已经学过牛顿定律：合力改变动量，物体在合力为零时保持静止或匀速直线运动。拿这套工具去看钥匙串，每一瞬间似乎都解释得通：拉下来时弹簧的拉力大于重力，合力向上，所以钥匙串向上加速；回到中间位置时合力为零；冲过中间位置后合力向下，所以它减速、停下、再回头……

可是这样逐瞬间地追下去，总感觉像在追着一辆永远追不上的车。我们更想要的是一张“全景地图”：这个系统整体上在干什么？它为什么不停在那个合力为零的位置——那里明明是所有力都平衡的地方？

直觉说：物体受力平衡了，就该安分了。实验说：恰恰相反，物体冲过平衡位置时速度最大，最不“安分”。它为什么不停下来？因为它有惯性——速度不会在力消失的瞬间自动变成零。冲过去之后，弹簧又被压缩（或拉伸），合力掉头往回拉，于是它减速、停下、折返。惯性让它“刹不住”，回复的拉力让它“逃不远”，两者一配合，来回振动就不可避免了。

这就是第一个重要的观念转变：振动不是一种需要新定律解释的怪现象，它就是“惯性”和“总想把自己拉回平衡位置的力”吵架的结果。哪里有这两样东西，哪里就有振动。而世界上几乎所有稳定的平衡位置附近，偏离一点点都会产生把你拉回去的力——碗底的小球、张紧的弦、压缩的气体、被弯曲的钢尺、甚至分子之间的化学键——所以世界才这么爱振动。严格地说，这是实验事实与数学模型的分工：“被扰动的系统常常来回振荡”是随处可见的实验事实；“回复力正比于偏离”是对平衡位置附近情况的数学近似；而“惯性与回复力的拉锯导致振荡”则是我们给出的物理解释。三者都不是“物体想回家”这样的目的论。

反过来想，什么东西不振动？一碗水静置在桌上，水面是平的；你用筷子点一下，水面先振几下才平静。一潭死水和一块果冻的区别不在“会不会动”，而在“推一下之后怎么回来”：果冻靠弹力把自己拉回原形，于是要振；水靠重力和表面张力回位，也会晃；而蜂蜜因为太黏，能量一下子就被耗散掉，只蠕回去不振。可见振动与否、振多久，是系统“倔强程度”“懒惰程度”和“耗散快慢”三方博弈的结果。这个视角比逐瞬间追力有用得多——它让你一眼看出一个陌生系统“大概会怎么动”。

\step{发明新概念}

现在像物理学家一样，把刚才那些口语一个个变成可以精确使用的概念。

\textbf{平衡位置}：系统不受扰动时安静待着的位置。对挂钥匙的弹簧，就是钥匙串静止下垂的位置；对秋千，是最低点；对钢尺，是尺子没弯时的形状。关键不在于这个位置本身，而在于：偏离它，力就要把你拉回去。还要注意平衡位置未必“什么都没有”：挂重物的弹簧在平衡位置处其实是拉长的，只是拉长产生的弹力恰好与重力抵消。我们关心的不是弹簧的绝对伸长，而是“离平衡位置还差多远”——把坐标原点挪到平衡位置，重力这位常驻嘉宾就被预先结算掉了，剩下的账全由回复力来记。这一步“换原点”的小手续，是让公式变干净的关键。

\textbf{回复力}：方向永远指向平衡位置、试图把系统拉回平衡位置的那个合力。注意，回复力不是一种和重力、弹力并列的新种类的力，它是合力扮演的角色：弹簧振子里由弹力和重力共同承担，单摆里由重力沿摆线方向的组分承担。这和“向心力”是同一个道理——名字描述的是任务，不是来源。

\textbf{简谐振动}：如果回复力的大小恰好正比于偏离平衡位置的距离，即
\[
  F = -kx,
\]
那么产生的振动叫简谐振动。这里 \(x\) 是偏离平衡位置的距离，\(k\) 是比例系数，弹簧的 \(k\) 就是中学里的劲度系数。负号是整个式子的灵魂：它保证力的方向永远和偏离的方向相反——你偏右，力朝左；你偏左，力朝右。弹簧在被拉伸和压缩不太过分时，非常精确地服从这个规律（胡克定律）。

\textbf{振幅、周期、频率、相位}：振动最远处离平衡位置的距离叫振幅 \(A\)，回答“振得多宽”；来回一次的时间叫周期 \(T\)，回答“一次多久”；单位时间内的次数叫频率 \(f=1/T\)，单位是赫兹（Hz，即每秒多少次）；相位则回答“此刻走到这一圈的哪一步了”——两个振幅和频率完全相同的秋千，如果一个在前一个在后，我们说它们相位不同。相位听起来抽象，其实你每天都在用：推秋千要“推在点上”，说的就是和秋千的相位配合。

有了这几个词，日常世界立刻变得可读。音叉上刻着“440 Hz”，意思是它每秒来回振动四百四十次，这正是音乐里的标准音；石英手表里是一片被精确切割的石英晶体，通上电后以 32768 Hz 的频率稳定振动——这个数是二的十五次方，电路只要连续减半十五次就得到恰好一秒一次的脉冲，手表的“滴答”就是这么来的；你手机里让整机嗡嗡作响的振动马达，不过是一个偏心小锤被电机带着高速旋转，相当于一个每秒上百次的强迫振动源；医院里的核磁共振成像，则是让人体内氢原子核在磁场中以特定频率“摇头晃脑”，再用同频率的无线电波去“推秋千”。从手表到人体，这些完全不同的系统共享同一套语言，因为它们满足同一个条件：都有平衡位置，都有回复力。

\begin{figure}[htbp]
\centering
\begin{tikzpicture}[thick,>=Latex]
  % 墙
  \draw[pattern=north east lines] (0,0) rectangle (0.25,2.4);
  \draw (0.25,0) -- (0.25,2.4);
  % 弹簧（螺旋示意）
  \draw[decorate,decoration={coil,aspect=0.5,segment length=4pt,amplitude=4pt}] (0.25,1.2) -- (3.3,1.2);
  % 滑块
  \draw[fill=blue!15] (3.3,0.8) rectangle (4.3,1.6);
  \node at (3.8,1.2) {$m$};
  % 地面
  \draw (2.8,0.6) -- (5.6,0.6);
  % 平衡位置线
  \draw[dashed,gray] (3.8,0.2) -- (3.8,2.4);
  \node[gray,below] at (3.8,0.2) {平衡位置};
  % 位移箭头
  \draw[<->,red] (3.8,2.15) -- (5.0,2.15) node[midway,above]{$x$};
  % 回复力箭头
  \draw[->,red,very thick] (4.3,1.2) -- (3.3,1.2) node[midway,below]{$F=-kx$};
\end{tikzpicture}
\caption{弹簧振子。偏离平衡位置 \(x\) 时，合力指向平衡位置；偏离越远，拉力越大。}
\end{figure}

\step{一句话模型}

\begin{oneline}
任何系统，只要偏离平衡位置时受到的合力总是把它往回拉、且拉力大致正比于偏离的距离，它就会在惯性（让它刹不住）和回复力（让它逃不远）的拉锯中以固定节奏来回振动；这个节奏只由系统自己的“倔强程度”（回复力系数 \(k\)）和“懒惰程度”（质量 \(m\)）决定，与你推得多狠无关。
\end{oneline}

把这句话拆开检查一遍。“惯性让它刹不住”解释了为什么冲过平衡位置；“回复力让它逃不远”解释了为什么还会回来；“节奏由系统决定”解释了钢尺实验里“扳得狠只改变响度、不改变音调”。“倔强”与“懒惰”是比喻：它像在描述弹簧硬不硬、物体沉不沉，不像的地方在于——这两个词在公式里有精确的含义，\(k\) 有单位、可以测量，不是形容词。

还要补一句“它不像什么”。这个模型不像“万物皆有灵性、都想回家”——回复力不是物体的愿望，而是结构决定的力学事实：弹簧被拉长，原子间距偏离了最省能的位置，力就出现了。它也不像“振动总会自动停下来才算正常”——恰恰相反，理想模型里振动永不停止，停下来是因为有能量漏出，这需要额外解释而不是默认设定。把“默认停止”当成常识，是亚里士多德时代留下的旧直觉，我们在讲运动学时已经拆过一次，这里再拆一次。

\step{数学翻译器}

直觉说：弹簧越硬，拉得越急，振动越快；物体越重，反应越迟钝，振动越慢。把这两句翻译成正比与反比关系，简谐振动的节奏就是
\[
  \omega=\sqrt{\frac{k}{m}},\qquad
  T=2\pi\sqrt{\frac{m}{k}},\qquad
  f=\frac{1}{2\pi}\sqrt{\frac{k}{m}}.
\]
其中 \(\omega\) 叫角频率，单位是弧度每秒，它和频率只差一个 \(2\pi\)：\(\omega=2\pi f\)。为什么出现平方根而不是简单正比，主线层可以这样理解：\(k\) 加倍，力加倍，加速度加倍，振动确实变快，但“快”的意思是更短时间内完成同样的往复，而加速度对时间的作用是“双重”的（速度是加速度积累出来的，位移又是速度积累出来的），两重积累叠加下来，时间只缩短为原来的 \(1/\sqrt{2}\)。严格的来源见数学桥层。

新公式到手，先用特殊情形检查它：

\begin{itemize}[leftmargin=2em]
  \item \(k\to 0\)（弹簧软到不存在）：\(T\to\infty\)，物体飘走不再回来。合理——没有回复力就没有振动。
  \item \(m\to\infty\)（挂上无限重的物体）：\(T\to\infty\)，振不起来。合理——惯性无穷大，谁也拉不动它。
  \item 公式里根本没有振幅 \(A\)：振幅加倍，周期不变。这正是饭勺实验里那个反直觉的事实，现在它从“怪事”变成了“预言”。注意它只在回复力严格正比于偏离时才成立，这会成为后面“模型边界”的主角。
\end{itemize}

单摆（秋千、挂着的饭勺）在小摆角时也近似是简谐振动。在给出公式之前，先用“量纲”这把尺子猜一猜它长什么样。周期是时间，单位是秒；单摆身上能和运动有关的参数只有摆长 \(L\)（米）、摆球质量 \(m\)（千克）和重力加速度 \(g\)（米每二次方秒）。要用这三样拼出一个“秒”，千克无处可去——除非质量根本不出现；而 \(L\) 除以 \(g\) 的单位是秒的平方，开个根号正好是秒。于是几乎唯一的可能就是 \(T\sim\sqrt{L/g}\)。这不是证明，只是“单位必须自洽”的硬约束，但它已经预言了两件事：周期与摆球质量无关，与摆长是平方根关系。用量纲先猜、再用实验或推导验证，是物理学家最常用的小抄之一。

严格的推导（见数学桥层）给出
\[
  T=2\pi\sqrt{\frac{L}{g}},
\]
\(L\) 是摆长，\(g\) 是重力加速度。检查：\(g\to 0\)（太空站里），\(T\to\infty\)——单摆彻底罢工，这将在脑洞实验里用到；\(L\) 变成四倍，周期只加倍，所以把秋千绳加长一倍并不会让摆动慢一倍，只慢大约四成。注意摆球的质量 \(m\) 根本不出现——这就是“先猜一猜”第一题的答案：体重加倍，周期几乎不变。

\textbf{一次带真实数据的估算。} 取一根摆长 \(L=1.0\) 米的单摆，地球表面 \(g\approx9.8\) 米每二次方秒：
\[
  T=2\pi\sqrt{\frac{1.0}{9.8}}\approx 2.0\ \text{秒}.
\]
来回一次正好约两秒——这不是巧合，而是钟表史上著名的“秒摆”：摆长约一米的钟摆，每摆动半周期滴答一秒。老座钟的钟摆长度就是照这个数定做的。反过来，用尺子量摆长、用手机计时测周期，你就能在家里测出当地的重力加速度，误差可以小到百分之几。振动把“时间”和“力”连在了一起，这正是它被用来做钟的原因。

\textbf{再来一次真实数据的估算：汽车悬挂。} 一辆普通轿车重约 1500 千克，四个车轮各分担约 375 千克。人坐进车里，车身会下沉大约 2 厘米，由 \(k=mg/x\) 可估出每个悬挂弹簧的劲度系数约为 \(375\times9.8/0.02\approx1.8\times10^{5}\) 牛每米。代入周期公式：
\[
  T=2\pi\sqrt{\frac{375}{1.8\times10^{5}}}\approx 0.3\ \text{秒},
\]
即车身上下颠簸的固有频率约 3 赫兹上下（实际设计得更低些，约 1 到 2 赫兹）。这解释了为什么碾过减速带后车身会以“一秒一两下”的节奏晃——这个节奏由车和弹簧决定，和你开车多快无关。也解释了为什么晕车的人往往对低频晃动最敏感：人腹腔内的脏器本身也是一组振子，固有频率恰好也在这个范围附近，外来的驱动若对上了节奏，五脏六腑都在共振，自然不会舒服。

再看能量。把振子拉到振幅处松手：那一刻速度为零，能量全部以弹性势能（或重力势能）的形式存着，大小为 \(\frac12 kA^{2}\)；经过平衡位置时弹簧恢复原状，势能清零，能量全部变成动能 \(\frac12 mv^{2}\)，此刻速度最大；冲到另一边又全部换回势能。振动的全过程，就是能量在“存起来”和“动起来”两种形态之间反复兑换，而总账户 \(\frac12 kA^{2}\) 保持不变。振幅加一倍，能量变四倍——所以把秋千推高一点点，要补充的能量远比看起来多。

能量视角还顺手回答了慢动作录像里那个观察：为什么振子在中间“一晃而过”、在两端“磨蹭”？因为总账户固定，势能少的地方动能必然多——平衡位置势能最低，速度就最大；端点势能全满，速度只能归零。这也解释了推秋千为什么在最低点推最划算：那里速度最大，同样大小的推力在同样作用距离内能传递最多的能量；在最高点附近推，秋千本来就快停了，推了半天也存不进多少。能量账本是比受力图更省事的分析工具，本书后面还会反复用到它。

\begin{mathbridge}
牛顿第二定律用于弹簧振子，给出
\[
  m\frac{d^{2}x}{dt^{2}}=-kx
  \quad\Longrightarrow\quad
  \frac{d^{2}x}{dt^{2}}+\omega^{2}x=0,
  \qquad \omega^{2}=\frac{k}{m}.
\]
这是一个二阶微分方程，意思是“找一种函数，它的二阶导数正比于自己的相反数”。正弦和余弦恰好满足：设 \(x(t)=A\cos(\omega t+\varphi)\)，求两次导数得 \(x''=-\omega^{2}x\)，自动成立。于是平方根的来源清楚了：它来自求导时角频率要“掉下来”两次。式中振幅 \(A\) 和初相位 \(\varphi\) 不由方程决定，而由初始条件（你一开始拉多远、松手时给不给初速度）决定——这从数学上解释了“节奏由系统决定，幅度由你决定”。单摆的严格方程是 \(\theta''=-(g/L)\sin\theta\)，它不是简谐的；只有当摆角很小、可以用 \(\sin\theta\approx\theta\) 替换时才变成简谐方程，这就是小摆角近似的出处。能量方面，\(\frac12 mv^{2}+\frac12 kx^{2}=\frac12 kA^{2}=\text{常量}\)，对无阻尼理想振子任何时刻都守恒。
\end{mathbridge}

\begin{figure}[htbp]
\centering
\begin{tikzpicture}[thick,>=Latex,scale=1.0]
  \draw[->] (0,0) -- (7.2,0) node[right]{$t$};
  \draw[->] (0,-1.6) -- (0,1.9) node[above]{$x$};
  \draw[blue,thick,domain=0:6.9,samples=100] plot (\x,{1.2*cos(\x*180)});
  \draw[dashed,gray] (0,1.2) -- (6.9,1.2);
  \draw[dashed,gray] (0,-1.2) -- (6.9,-1.2);
  \draw[<->,red] (6.45,0) -- (6.45,1.2) node[midway,right]{$A$};
  \draw[<->,red] (1.0,-1.5) -- (3.0,-1.5) node[midway,below]{$T$};
  \fill (1.0,1.2) circle (1.2pt); \fill (3.0,1.2) circle (1.2pt);
\end{tikzpicture}
\caption{简谐振动的位移—时间图像是一条正弦（或余弦）曲线：相邻两个波峰相隔一个周期 \(T\)，峰到平衡位置的距离是振幅 \(A\)。}
\end{figure}

\step{模型的边界}

简谐模型干净漂亮，但它有三个明确的适用条件，越界就会出错。

\textbf{条件一：回复力要正比于偏离。} 弹簧拉得太过分会永久变形，胡克定律失效；单摆的 \(\sin\theta\approx\theta\) 只在小角度成立，摆角大到三四十度以上，周期会明显变长（摆角 \(60^{\circ}\) 时周期约比小角近似长百分之七）。“振幅不影响周期”从此降为近似结论。这是直觉最容易出错的地方：日常摆角小，误差小，于是我们把近似当成了定律。

\textbf{条件二：没有能量漏出。} 真实振子总有摩擦和空气阻力，能量不断变成热散掉，振幅随时间衰减——这叫阻尼。钥匙串最终停下，不是谁的“力气用完了”，而是机械能转移成了内能。阻尼并不总是坏事：汽车减震器（避震筒）就是故意把弹簧振子的能量“抽走”，否则碾过一个坑你要晃到天黑。这里也修正一个流行误解：不是弹簧让车舒服，是“弹簧负责缓冲、阻尼器负责收尾”这对搭档让车舒服；只有弹簧没有阻尼的车，才真叫颠。

阻尼的三档“火候”值得分清。阻尼太小，振子要来回很多圈才停，像松开的吉他弦；阻尼恰好到某个临界值，振子一次就回到平衡位置、不再过头，这叫临界阻尼——好的车门、好的电流表指针都调在这一档，追求的是“回位快且不晃”；阻尼过大，振子像掉进蜂蜜里，缓慢蠕回平衡位置，连一次都振不起来。同一个弹簧配上不同阻尼，性格完全不同。工程师调减震器，调的其实就是这档火候。

\textbf{条件三：别把共振当成万能解释。} 塔科马大桥常被教科书当作“风引起共振吹垮大桥”的标准案例，但更细致的事故分析表明，风中桥梁的扭振更接近一种“气动自激振动”（颤振）：桥的扭动改变了气流，气流反过来加剧扭动，即使风速恒定也能自我放大。它和推秋千有共同的灵魂——能量在正确的相位上持续输入——但机制比一般共振更微妙。把一切都归因于“共振”，是对共振的滥用；真正的教训是“振动系统在持续供能时可能失控”，这在抗风设计里是严肃的工程科目。

\begin{figure}[htbp]
\centering
\begin{tikzpicture}[thick,>=Latex,scale=1.0]
  \draw[->] (0,0) -- (7.0,0) node[right]{驱动频率};
  \draw[->] (0,0) -- (0,3.1) node[above]{振幅};
  \draw[blue,thick] plot[smooth] coordinates {(0,0.5) (1.5,0.7) (2.6,1.4) (3.2,2.8) (3.8,1.4) (5.0,0.8) (6.5,0.55)};
  \draw[red,thick] plot[smooth] coordinates {(0,0.45) (1.5,0.6) (2.6,1.0) (3.2,1.5) (3.8,1.0) (5.0,0.7) (6.5,0.5)};
  \draw[dashed,gray] (3.2,0) -- (3.2,2.8);
  \node[gray,below] at (3.2,0) {$f_0$};
  \node[blue] at (4.6,2.4) {阻尼小：尖峰};
  \node[red] at (5.6,1.2) {阻尼大：平缓};
\end{tikzpicture}
\caption{受迫振动的共振曲线：驱动频率接近固有频率 \(f_0\) 时振幅最大；阻尼越小，峰越尖越高。}
\end{figure}

\begin{challenge}
受迫振动的方程是 \(x''+2\beta x'+\omega_0^2 x=(F_0/m)\cos\omega t\)，其中 \(\beta\) 描述阻尼，\(\omega\) 是驱动频率。它的稳态解是同频率的受迫振动，振幅
\[
  A(\omega)=\frac{F_0/m}{\sqrt{(\omega_0^{2}-\omega^{2})^{2}+4\beta^{2}\omega^{2}}}.
\]
检查两个极限：\(\omega\to 0\) 时 \(A\to F_0/k\)，即慢慢推等于静力拉伸，合理；\(\beta\to 0\) 且 \(\omega\to\omega_0\) 时分母趋零，振幅发散——无阻尼共振在数学上振幅无限大，现实中则意味着“迟早有东西先坏掉”。工程上用品质因数 \(Q=\omega_0/(2\beta)\) 刻画共振峰的尖锐程度：音叉的 \(Q\) 可达上千，敲一下响很久；汽车悬挂被故意设计成 \(Q\) 接近一，一个坑只颠一下。同一个数学结构还统治着 RLC 电路和激光谐振腔——振动是唯一的，载体是多样的。
\end{challenge}

\step{回头解释开场现象}

现在回头批改开场的三个悬念。

\textbf{秋千为什么能被轻轻推高。} 秋千是一个近似单摆，有自己的固有节奏 \(T=2\pi\sqrt{L/g}\)，而且小摆角下这个节奏几乎不随摆幅改变。旁人的轻推本质上是一个小驱动：只要推的节奏对上秋千的固有节奏，每一次推都和秋千的运动同相位，每一圈都往系统里净存一点能量，振幅就一圈圈涨上去——这就是共振。推力小没关系，关键是“每一圈都存在同一张存折里”。反之乱推时，有时推在秋千回程上，那一推是从存折里取钱，振幅自然涨不起来。这同时回答了先猜一猜的第二题：关键在时机，不在力气。

\textbf{饭勺为什么宽窄同节奏。} 小摆角下单摆周期 \(T=2\pi\sqrt{L/g}\) 不含振幅：晃得宽时，摆球走的路程长，但重力沿摆线的分量也同比例变大，两种效果恰好抵消，周期不变。这个性质叫等时性，伽利略正是靠它（以及自己的脉搏当表）想到用摆来计时。你的直觉之所以猜错，是因为直觉把“走得远”直接翻译成“花的时间长”，忘了力也在同步变大。

\textbf{大桥为什么会振到散架。} 桥面不是绝对刚体，它有自己的固有频率。稳定的风绕过桥面形成交替脱落的涡旋，相当于一个持续的驱动；一旦驱动的节奏接近桥的某个固有频率，共振放大就开始了，再加上前述气动自激的推波助澜，振幅最终超过了钢结构的承受极限。此后所有悬索桥设计都要做风洞试验，让桥的固有频率躲开常见风况的驱动频率，并用阻尼结构吸收能量。同一原理好的一面是音乐：吉他的琴箱、小提琴的面板都被设计成在与琴弦相同的频率附近容易振动，弦的微小振动通过共振被放大成我们听得见的声音。共振既能制造音乐，也能摧毁桥梁，区别只在于我们想让能量存进哪里。

摩天大楼面对的问题和桥一模一样，解法也一脉相承。台北 101 大楼在八十八层附近悬挂着一颗直径五米半、重约 660 吨的钢球，用钢缆吊着，相当于一个巨型单摆。台风推动大楼以自身固有频率摇晃时，工程师故意把这颗球的摆动节奏调到大楼的固有频率上，让球与楼“对着干”：楼往东晃，球拖在后面，通过连接两者的阻尼油缸把楼的振动能量持续抽走变成热。这个装置叫调谐质量阻尼器，本质是把“共振”这位破坏王招安成保镖——既然能量总要有个去处，不如预先给它安排一个安全的去处。地震高发的日本、风力强劲的沿海地区，许多超高层建筑体内都藏着类似的“镇楼之球”。

\step{脑洞实验与挑战题}

\begin{thought}[在空间站里称体重]
空间站里一切失重，站上体重秤读数为零——不是你没质量，而是秤的原理（压弹簧直到支持力平衡重力）失效了。那宇航员怎么监测自己的体重？答案是利用本章的公式：让宇航员坐在一把连着已知弹簧的椅子上，拉离平衡位置松手，测振动周期 \(T\)，代入 \(m=kT^{2}/(4\pi^{2})\) 直接算出质量。空间站里确实装着这样的装置。注意这个方法的深刻之处：它测量的不是“重力下的重量”，而是“惯性有多大”——质量的这个含义（惯性质量）在失重环境里依然完好。单摆在这里会彻底罢工（\(g\to0\)，\(T\to\infty\)），弹簧振子却照常工作，两种振动对重力的依赖完全不同。再进一步想：如果把这套弹簧椅送到遥远的星际空间，测量结果会变吗？
\end{thought}

振动还是通往本书后文的桥。一根吉他弦可以看成无数个小振子手拉手排成一排，一个振子动起来会牵着邻居动——一群振子的集体舞蹈就是波，那是第 17 章的主题。复杂杂乱的振动又能拆成许多简单频率的叠加，那是第 18 章傅里叶思想的入口。再往后，晶格中原子被化学键这根“弹簧”拴在平衡位置附近，它们的热振动与物体的温度密切相关；而把量子力学用到谐振子上，会发现振子的能量竟是一份一份的，而且永远不可能完全静止——零点能的故事，在量子部分等着你。甚至电路也不例外：电容和电感组成的回路里，电场的能量和磁场的能量像动能和势能一样来回兑换，那也是一台不折不扣的振子，收音机靠它选台。学会看这一个弹簧，你就拿到了一大半物理学的钥匙。

\begin{puzzles}
\item 同一根弹簧挂同一块重物：在地球表面上下振动的周期，与搬到月球表面（重力约为地球的六分之一）后相比，会怎样变化？若换成同摆长的单摆做同样搬迁呢？（提示：检查两个周期公式里有没有 \(g\)；弹簧的 \(k\) 由弹簧自身决定。）
\item 荡秋千时，高手会在最低点“蹲下”、在最高点“站起”，不靠任何人推就能越荡越高。用能量观点解释他从哪里“偷”来了能量。（提示：站起时身体的重心升高，而那一刻他在对抗把自己往下压的力做功；蹲下发生在速度最大的位置，代价小。这相当于每个周期改变一次有效摆长，属于参数驱动，和普通推力驱动是两种供能方式。）
\item 老式收音机转动选台旋钮时，内部是在调节一个电学“振子”（电容和电感组成的回路）的固有频率。用这个模型说明：为什么收到一个台时其他台都安静了，以及为什么有的收音机两个相邻电台会“串台”。（提示：对照共振曲线的形状，串台对应峰太宽——即阻尼太大、选择性太差。）
\item 一座设计不佳的桥，行人步频恰好接近桥的固有频率时会越晃越厉害（现实中真发生过，伦敦千禧桥开通两天就因晃动关闭整改）。整改方案不是加粗钢梁，而是加装阻尼器。为什么“让桥更结实”不如“让桥更耗能”？（提示：共振时振幅由输入功率和耗散功率的平衡决定；把 \(Q\) 值压下来，峰就塌了。）
\end{puzzles}
